% =====================================================================
% Stage 5 -- Test Case Specifications.
%
% The final stage of the refinement flow. Every verifiable artifact in the
% preceding stages carries a verify obligation -- \reqcover{verify}{TC}, or a
% \reqexpect naming the specific TC; each is discharged here by a test-case
% specification that \verify's it.
%
% These specifications are drafts -- the objective and the shape of each
% test are fixed, but the concrete fixtures, sample responses, and exact
% assertions are left to be filled in when the tests are implemented (an
% IMP artifact \satisfy-ing each TC).
% =====================================================================

We assume the compliance tests will be driven by
\href{https://github.com/http-tests/cache-tests}{\texttt{http-tests/cache-tests}},
an existing, actively maintained test suite for HTTP caches (Appendix~\ref{app:cache-tests}
details it fully). It ships its own controllable origin server, so the ``mocked
upstream'' needed for the test scenarios does not need to be built: each test
case addresses the origin with a unique ID and receives back exactly the
headers, timing, and body that scenario requires.

What remains to be built is the \emph{test subject}: a minimal Pingora proxy
that enables \texttt{pingora-cache} with a single upstream peer pointed at the
cache-tests origin and the file-system backend as its only \texttt{Storage}
implementation. This proxy is a disposable test fixture, not a project
deliverable --- it exists solely to host the storage backend under test for
the duration of a compliance run and ships nothing.\footnote{Running it also
introduces a Node.js/npm dependency into the test or CI environment, since the
cache-tests CLI client is JavaScript-based.} \reqref{TC-COMPLIANCE.RFC-01}
runs the cache-tests CLI against that proxy twice --- once with the stock
\texttt{MemCache} as its \texttt{Storage} and once with the file-system
backend --- and compares the two result sets test by test. Because
\texttt{pingora-cache} does not itself claim full RFC~9111 conformance, the
stock run is the reference: a test that fails on both runs is a limitation of
\texttt{pingora-cache}, not of this backend, which is exactly what
\reqref{HLR-COMPLIANCE.RFC-01} requires to be distinguished
(Appendix~\ref{app:cache-tests}).

% ---------------------------------------------------------------------
% Standards compliance.
% ---------------------------------------------------------------------
\begin{requirement}{TC-COMPLIANCE.RFC-01}[draft][verify]{No regression against MemCache}

  \localreqdiagram*{tc01}{\reqid}{\reqdesc}

  \textbf{Objective:} Show that the file-system backend adds no RFC 9111/9110
  deviation to the proxy's observable caching behaviour, by comparing it with
  the same proxy using the stock \texttt{MemCache}.

  \textbf{Setup:} One proxy build whose \texttt{Storage} implementation is
  selectable at start --- the stock \texttt{MemCache} or the file-system
  backend --- with its upstream peer pointed at the \texttt{cache-tests}
  origin server, and all other configuration identical. \texttt{cache-tests}
  supplies its own corpus of per-test-case responses and expected outcomes;
  Appendix~\ref{app:cache-tests} (\S~\ref{app:cache-tests-mapping}) lists
  which \texttt{cache-tests} files exercise each \texttt{FRE-COMPLIANCE.*}
  requirement, and which requirements it does not reach.

  \textbf{Steps:}
  \begin{enumerate}
    \item Start the proxy with \texttt{MemCache}; run the \texttt{cache-tests}
      CLI against it, exporting the results as JSON.
    \item Start the proxy with the file-system backend on an empty cache
      directory; run the CLI again the same way.
    \item Compare the two result sets test by test, ranking outcomes pass,
      then suboptimal, then fail. Re-run against both backends every test
      whose outcome differs, to separate a reproducible difference from a
      timing-dependent one.
  \end{enumerate}

  \textbf{Expected result:} No test has a reproducibly lower-ranked outcome
  with the file-system backend than with \texttt{MemCache}. Tests that fail on
  both runs are recorded as \texttt{pingora-cache} limitations, outside this
  requirement. Observation-only results that differ are listed for review.
  (Restart scenarios are outside this test --- \texttt{cache-tests} cannot
  restart the process under test --- and are covered by the per-requirement
  test cases.)

  \reqdiaglink{overview}{\textit{View the domain overview.}}

  \reqcover{satisfy}{IMP}

  \verify{HLR-COMPLIANCE.RFC-01}
\end{requirement}

\begin{requirement}{TC-COMPLIANCE.HEADERS-01}[draft][verify]{Header fields round-trip faithfully}

  \localreqdiagram*{tc02}{\reqid}{\reqdesc}

  \textbf{Objective:} A stored response returns the exact response head it was
  given --- status code, version, and every header field.

  \textbf{Setup:} The backend driven directly through the \texttt{Storage}
  interface by a test harness; a ``restart'' constructs a new backend instance
  over the same cache directory.

  \textbf{Steps:}
  \begin{enumerate}
    \item Store a response whose head exercises each shape fidelity must not
      depend on: repeated field names in a deliberate order, an unrecognised
      field, an empty-valued field, a field value carrying opaque non-ASCII
      octets (obs-text), a \texttt{Set-Cookie}-style field repeated across
      lines, a non-\texttt{200} status code, header sets of eight consecutive
      serialised sizes stored with no metadata headroom, so that each
      metadata region ends exactly where the body begins
      (\reqref{SDS-LAYOUT.OBJECTFILE-01}), and a header set many times the
      default headroom.
    \item Look the object up; then restart the process and look it up again.
  \end{enumerate}

  \textbf{Expected result:} Both lookups return the status code and version
  unchanged and every field with its name and value intact byte for byte, the
  relative order of same-named fields preserved; nothing is dropped, added,
  reordered, or normalised, regardless of the size of the header set. Where
  the metadata region is exactly full, the body that follows it is intact.

  \reqdiaglink{overview}{\textit{View the domain overview.}}

  \reqcover{satisfy}{IMP}

  \verify{FRE-COMPLIANCE.HEADERS-01}
\end{requirement}

\begin{requirement}{TC-COMPLIANCE.TRAILERS-01}[draft][verify]{Trailer fields kept distinct}

  \localreqdiagram*{tc03}{\reqid}{\reqdesc}

  \textbf{Objective:} Show, by inspection, that trailer fields are never
  returned as header fields. No test can present a trailer through the
  targeted interface (\reqref{FRE-COMPLIANCE.TRAILERS-01}), so this
  specification records an inspection procedure rather than a test.

  \textbf{Steps (inspection):}
  \begin{enumerate}
    \item Inspect the \texttt{Storage} and \texttt{HandleMiss} signatures of
      the targeted \texttt{pingora-cache} version: confirm that a store
      receives header metadata only as a \texttt{CacheMeta} and body bytes
      only through \texttt{write\_body}, with no trailer parameter.
    \item Inspect the on-disk record format
      (\reqref{SDS-LAYOUT.OBJECTFILE-01}) and its implementation: confirm
      that the stored header segment is written only from the
      \texttt{CacheMeta} handed to the store or metadata update, and that no
      other section could supply fields to it.
  \end{enumerate}

  \textbf{Expected result:} The interface offers no trailer channel, and
  the record format has no path by which a trailer-originated field could
  enter the stored header segment. Repeat the inspection whenever the
  targeted \texttt{pingora-cache} version changes.

  \reqdiaglink{overview}{\textit{View the domain overview.}}

  \reqcover{satisfy}{IMP}

  \verify{FRE-COMPLIANCE.TRAILERS-01}
\end{requirement}

\begin{requirement}{TC-COMPLIANCE.BODY-01}[draft][verify]{Body round-trips byte for byte}

  \localreqdiagram*{tc04}{\reqid}{\reqdesc}

  \textbf{Objective:} The served body equals the stored body exactly, for full
  and for seeked (partial) reads.

  \textbf{Setup:} The backend driven directly through the \texttt{Storage}
  interface by a test harness; a ``restart'' constructs a new backend instance
  over the same cache directory.

  \textbf{Steps:}
  \begin{enumerate}
    \item Store responses with random binary bodies at boundary sizes: empty,
      one byte, exactly at and one byte either side of the backend's
      write-chunk and file-system block sizes, and a large multi-chunk body.
    \item Retrieve each, including after a restart, and compare to the bytes
      written.
    \item Where the hit handler advertises seek support, read sub-ranges of the
      large body --- the start, an interior slice spanning a chunk boundary,
      and the tail --- and compare each against the corresponding slice of the
      stored bytes.
  \end{enumerate}

  \textbf{Expected result:} Each full read is identical byte for byte and in
  order --- no truncation, padding, or re-encoding; each seeked read delivers
  exactly the requested slice, so range responses built from the reader carry
  precisely the stored bytes.

  \reqdiaglink{overview}{\textit{View the domain overview.}}

  \reqcover{satisfy}{IMP}

  \verify{FRE-COMPLIANCE.BODY-01}
\end{requirement}

\begin{requirement}{TC-COMPLIANCE.FRESHNESS-01}[draft][verify]{Freshness and age inputs preserved}

  \localreqdiagram*{tc05}{\reqid}{\reqdesc}

  \textbf{Objective:} Freshness and age computed from a retrieved object match
  those computed before it was stored.

  \textbf{Setup:} The backend driven directly through the \texttt{Storage}
  interface by a test harness; a ``restart'' constructs a new backend instance
  over the same cache directory, with freshness and age evaluated by
  the functions of \texttt{pingora-cache} itself over the returned \texttt{CacheMeta}.

  \textbf{Steps:}
  \begin{enumerate}
    \item Store responses exercising \texttt{Date}, \texttt{Expires}, \texttt{Age},
      and \texttt{max-age}/\texttt{s-maxage}, including boundary and degenerate
      values: \texttt{Age: 0}, delta-seconds at the overflow cap
      (\texttt{2147483648}), \texttt{max-age=0}, an already-past
      \texttt{Expires}, and a syntactically invalid \texttt{Expires}.
    \item Evaluate freshness lifetime and current age before storing and after
      retrieving (including across a restart), using the same clock reference.
    \item Apply a metadata update carrying new freshness fields
      (\reqref{FRE-COMPLIANCE.REVALIDATION-01}) and re-evaluate from a fresh
      lookup.
  \end{enumerate}

  \textbf{Expected result:} The pre-store and post-retrieve evaluations agree,
  including the internal created/updated instants of \texttt{pingora-cache}; boundary and invalid
  values round-trip as opaque bytes, un-normalised; after the update, evaluation
  reflects the updated field set and the new updated-at instant.

  \reqdiaglink{overview}{\textit{View the domain overview.}}

  \reqcover{satisfy}{IMP}

  \verify{FRE-COMPLIANCE.FRESHNESS-01}
\end{requirement}

\begin{requirement}{TC-COMPLIANCE.VALIDATORS-01}[draft][verify]{Validator fields preserved}

  \localreqdiagram*{tc06}{\reqid}{\reqdesc}

  \textbf{Objective:} \texttt{ETag} and \texttt{Last-Modified} survive storage
  unchanged, including the weak indicator.

  \textbf{Setup:} The backend driven directly through the \texttt{Storage}
  interface by a test harness; a ``restart'' constructs a new backend instance
  over the same cache directory.

  \textbf{Steps:}
  \begin{enumerate}
    \item Store responses with a strong \texttt{ETag}, a weak (\texttt{W/})
      \texttt{ETag}, an empty entity tag (\texttt{ETag: ""}), a
      \texttt{Last-Modified} date, and one response carrying both validators.
    \item Retrieve and compare the validator fields to those stored, including
      after a restart.
    \item Apply a metadata update that changes the validators
      (\reqref{FRE-COMPLIANCE.REVALIDATION-01}); retrieve and compare again.
  \end{enumerate}

  \textbf{Expected result:} The validators are returned byte-exactly, with the
  \texttt{W/} marker retained and the empty tag intact, so conditional
  revalidation can be issued and matched; after the update, every lookup
  returns the updated validators.

  \reqdiaglink{overview}{\textit{View the domain overview.}}

  \reqcover{satisfy}{IMP}

  \verify{FRE-COMPLIANCE.VALIDATORS-01}
  \verify{FRE-COMPLIANCE.REVALIDATION-01}
\end{requirement}

\begin{requirement}{TC-COMPLIANCE.REVALIDATION-01}[draft][verify]{Metadata freshened atomically}

  \localreqdiagram*{tc07}{\reqid}{\reqdesc}

  \textbf{Objective:} A metadata update after a \texttt{304} replaces metadata
  while leaving the body intact, applies only to the object it validated, is
  observed atomically and is visible to lookups when \texttt{update\_meta}
  returns; an update that outgrows the record's reserved space withdraws the
  object, observably, so that the next request fetches it again.

  \textbf{Setup:} The backend driven directly through the \texttt{Storage}
  interface, with a test hook able to make a withdrawal's header write fail,
  and log and metric capture.

  \textbf{Steps:}
  \begin{enumerate}
    \item Store an object, then apply a metadata update with a newer header set;
      repeat with an update whose header set is large enough that it no longer
      fits the space the layout reserved ahead of the body
      (\reqref{FDS-STORAGE.METAUPDATE-01}), while a client reads the object
      through a hit opened before that update.
    \item Run lookups concurrently with each update.
    \item As each \texttt{update\_meta} call returns, immediately look the
      object up, before any other task runs, and inspect the record on disk.
    \item Restart the process and look the object up again.
    \item Apply a metadata update to a key with no stored object.
    \item Configure zero reserved headroom and repeat the small-update case
      from step 1.
    \item After an update has withdrawn a record, store the object again, as
      the next request's fetch does, and look it up.
    \item Apply sixteen updates that each outgrow the reserved space, all at
      once. With a store of the key from the origin in flight, apply an
      update that outgrows the record, then let the store finish.
    \item Make the backend read-only (\reqref{FDS-COMPAT.READONLY-01}) and
      apply an update that fits and one that does not.
    \item Make the header write of a withdrawal fail; then apply the update
      again.
    \item Read the counter
      \texttt{pingora\_durable\_meta\_headroom\_exceeded\_total} once the
      backend is built. Then, for every withdrawal above, and for the one
      whose header write failed, read the captured log and the counter.
    \item In a build without the backend's \texttt{metrics} feature, apply one
      update that withdraws the record; read the captured log and the
      registry.
    \item Store an object, then a newer object of the same key, created later;
      then apply the update a \texttt{304} for the first object brings: new
      header fields under the first object's creation time.
  \end{enumerate}

  \textbf{Expected result:} The body is unchanged and still retrievable under the
  same key; every concurrent lookup sees either the pre-update or post-update
  metadata, never a mixture, and never a body mismatched to its metadata;
  when each call that fits returns, a lookup answers with the update and the
  record on disk holds it; after the restart, lookups return the updated
  metadata, never the superseded set; the missing-key update reports an
  error and materialises no object; the update that outgrows the reserved
  space reports \texttt{false} and withdraws the record in place, the file
  keeping its length: the client already reading it gets the whole body, and
  every later lookup misses, after the restart too; with zero headroom even
  the small update withdraws the record; the store made after a withdrawal
  replaces the withdrawn file, with a region sized for its own header set,
  and lookups hit it; of the sixteen updates, one withdraws the record and
  the others report that the key holds no object; the update made during
  the store from the origin withdraws the record, and that store then lands
  and replaces it; the updates once the backend is read-only fail with an
  error and change nothing; the withdrawal whose header write fails reports
  an error and leaves the record as it was, served with its old metadata,
  and the next attempt withdraws it; the counter reads zero once the
  backend is built; each withdrawal logs one warning naming
  the key's hash and the headroom shortfall, and raises the counter by one,
  while the failed one is neither logged nor counted; in the build without
  the \texttt{metrics} feature the warning is still logged and the counter
  is not registered; and the update for the replaced object is reported not
  applied, and the newer object stands whole.

  \reqdiaglink{overview}{\textit{View the domain overview.}}

  \reqcover{satisfy}{IMP}

  \verify{FRE-COMPLIANCE.REVALIDATION-01}
  \verify{FRE-COMPLIANCE.METAATOMIC-01}
  \verify{FRE-COMPLIANCE.METAPERSIST-01}
  \verify{FRE-COMPLIANCE.METAABSENT-01}
  \verify{FDS-STORAGE.METAUPDATE-01}
  \verify{FDS-OBSERVABILITY.HEADROOM-01}
\end{requirement}

\begin{requirement}{TC-COMPLIANCE.VARY-01}[draft][verify]{Vary variants stored and selected}

  \localreqdiagram*{tc08}{\reqid}{\reqdesc}

  \textbf{Objective:} Multiple variants of one primary key are stored and
  correctly selected.

  \textbf{Setup:} The backend behind a \texttt{pingora-cache}
  \texttt{HttpCache} driven by a test harness, so that lookups, variant
  selection, and purges follow the framework's own call sequence; a
  ``restart'' constructs new instances over the same cache directory.

  \textbf{Steps:}
  \begin{enumerate}
    \item Store two responses under the same primary key with different variance
      keys derived from \texttt{Vary}.
    \item Look up with requests matching each variant, and with one matching
      neither.
    \item Replace one variant with a newer response and apply a metadata update
      to it; look both variants up again. Then purge that variant's own
      variance-augmented key through the backend's purge operation, and look
      the other variant up again.
    \item Restart the process and look each variant up once more.
  \end{enumerate}

  \textbf{Expected result:} Each matching lookup returns its own variant; the
  non-matching lookup reports a miss; one variant is never served for another;
  mutating one variant leaves the other byte-identical; the full variant set
  survives the restart, each variant still retrievable under its
  variance-augmented key.

  \reqdiaglink{overview}{\textit{View the domain overview.}}

  \reqcover{satisfy}{IMP}

  \verify{FRE-COMPLIANCE.VARY-01}
  \verify{FRE-COMPLIANCE.VARYISOLATION-01}
\end{requirement}

\begin{requirement}{TC-COMPLIANCE.PURGE-01}[draft][verify]{Purge makes all variants unreachable}

  \localreqdiagram*{tc09}{\reqid}{\reqdesc}

  \textbf{Objective:} After a purge of a key, neither the response nor any of its
  variants is ever served again.

  \textbf{Setup:} The backend behind a \texttt{pingora-cache}
  \texttt{HttpCache} driven by a test harness, so that lookups, variant
  selection, and purges follow the framework's own call sequence; a
  ``restart'' constructs new instances over the same cache directory.

  \textbf{Steps:}
  \begin{enumerate}
    \item Store several \texttt{Vary} variants under one key.
    \item Issue a purge for that key (the proxy's purge path addresses the primary
      slot), then look the primary and each variant up again, through
      the primary-slot-first lookup in \texttt{pingora-cache}.
    \item Purge the same key a second time, and purge a key that was never
      stored.
    \item Restart the process, once cleanly and once after killing it, and
      look the primary and each variant up again. (Purge durability across a
      host crash is exercised by \reqref{TC-PERSISTENCE.CRASHSAFE-01}.)
    \item Repeat the purge while a reader is mid-stream on one variant and
      while a same-key store is in flight; let both run to completion.
    \item Restart the process once more, letting Eviction Recovery's scan run
      (\reqref{FDS-EVICTION.STARTUP-01}), and inspect the cache directory and
      the eviction manager for the variant records.
  \end{enumerate}

  \textbf{Expected result:} Every lookup --- primary and variants --- reports a
  miss after the purge; an invalidated response is never served; the repeated
  and never-stored purges are harmless no-ops that report nothing found; no
  purged object reappears after the restart or crash; the mid-stream reader
  runs to completion or fails cleanly --- never another object's bytes --- and
  its storage is not reused while referenced; the racing store publishes
  nothing, since the purge ended it; after the final restart, the variant
  records the purge orphaned are still in the cache directory, each admitted
  to the eviction manager with its record size, so that eviction removes them
  like any other object.

  \reqdiaglink{overview}{\textit{View the domain overview.}}

  \reqcover{satisfy}{IMP}

  \verify{FRE-COMPLIANCE.PURGE-01}
  \verify{FRE-COMPLIANCE.PURGEABSENT-01}
  \verify{FRE-COMPLIANCE.PURGEDURABLE-01}
  \verify{FRE-COMPLIANCE.READERISOLATION-01}
\end{requirement}

\begin{requirement}{TC-COMPLIANCE.INCOMPLETE-01}[draft][verify]{Incomplete responses never served}

  \localreqdiagram*{tc10}{\reqid}{\reqdesc}

  \textbf{Objective:} A store that does not run to completion is not later served
  as if complete.

  \textbf{Setup:} The backend behind a \texttt{pingora-cache}
  \texttt{HttpCache} driven by a test harness, so that lookups, variant
  selection, and purges follow the framework's own call sequence; a
  ``restart'' constructs new instances over the same cache directory.

  \textbf{Steps:}
  \begin{enumerate}
    \item Begin storing a response and, while the store is in flight, attach a
      follower streaming the partial hit
      (\reqref{FRE-COMPAT.STREAMING-01}, \reqref{FRE-COMPLIANCE.PARTIALABORT-01}).
    \item Abort the store before the body is fully written (simulating an
      upstream disconnect), then drop its miss handler; inspect the cache
      directory.
    \item Look the key up, including after a restart.
    \item Repeat, but instead of an orderly abort kill the process mid-write
      (no abort path runs); restart, look the key up, and inspect the cache
      directory for the per-writer temporary left behind.
  \end{enumerate}

  \textbf{Expected result:} Every lookup reports a miss --- the only observable
  the \texttt{Storage} interface permits for an incomplete entry --- and a
  partial body is never returned as a complete response; the follower's
  in-flight stream terminates with an error, never as a clean end of body; the
  orderly-aborted store's temporary is gone by the time its handler is dropped;
  the crash-orphaned temporary is never served or promoted, and is gone once
  the restarted process's restart-time scan has run. States a real host crash
  can leave behind are exercised by \reqref{TC-PERSISTENCE.CRASHSAFE-01}.

  \reqdiaglink{overview}{\textit{View the domain overview.}}

  \reqcover{satisfy}{IMP}

  \verify{FRE-COMPLIANCE.INCOMPLETE-01}
  \verify{FRE-COMPLIANCE.PARTIALABORT-01}
  \verify{FRE-COMPAT.STARTRECLAIM-01}
\end{requirement}

\begin{requirement}{TC-COMPLIANCE.REPLACE-01}[draft][verify]{Stores replace atomically}

  \localreqdiagram*{tc17}{\reqid}{\reqdesc}

  \textbf{Objective:} Storing a new response under an occupied key replaces the
  old object atomically; readers and racing writers are never torn.

  \textbf{Setup:} The backend driven directly through the \texttt{Storage}
  interface by a test harness; a ``restart'' constructs a new backend instance
  over the same cache directory.

  \textbf{Steps:}
  \begin{enumerate}
    \item Store object A under a key and attach a reader mid-stream on A's body.
    \item Store object B --- different metadata and a different, larger body ---
      under the same key; run lookups while B's store is in flight, then let it
      complete and look up again, including after a restart.
    \item Drain the reader attached in step 1 and compare what it received.
    \item Race two same-key stores (C and D) to completion and look up.
  \end{enumerate}

  \textbf{Expected result:} In-process lookups during B's store return a
  partial hit of B (\reqref{FRE-COMPAT.STREAMING-01}), and a lookup that
  cannot see B's in-flight write returns A whole; after B's
  store completes, every lookup --- also after the restart --- returns B in its
  entirety, never A's metadata with B's body or any other mixture; the reader
  of A receives exactly A's bytes to a clean end of body (or a clean error),
  never bytes of B, and A's space is reclaimed only once no reader references
  it; after the race, lookups return exactly one winner (the last published),
  intact.

  \reqdiaglink{overview}{\textit{View the domain overview.}}

  \reqcover{satisfy}{IMP}

  \verify{FRE-COMPLIANCE.REPLACE-01}
  \verify{FRE-COMPLIANCE.READERISOLATION-01}
\end{requirement}

% ---------------------------------------------------------------------
% Persistence.
% ---------------------------------------------------------------------
\begin{requirement}{TC-PERSISTENCE.CRASHSAFE-01}[draft][verify]{Host crash serves no bad content}

  \localreqdiagram*{tc27}{\reqid}{\reqdesc}

  \textbf{Objective:} In every state a host crash or power loss could leave on
  disk, the restarted cache serves no object whose content differs from what
  was stored and no object whose purge it had acknowledged; objects that did
  not survive are misses.

  \textbf{Setup:} The cache directory on an ext4 file system whose block device
  is a device-mapper \texttt{log-writes} target, which logs every write with its
  flush and FUA ordering so that the device can be replayed to any point a power
  failure could have left it;\footnote{\href{https://docs.kernel.org/admin-guide/device-mapper/log-writes.html}{Linux
  kernel documentation: ``dm-log-writes''}, with its replay tool at
  \url{https://github.com/josefbacik/log-writes}.} the backend driven directly
  through the \texttt{Storage} interface, with the log marked after each
  operation the backend acknowledges.

  \textbf{Steps:}
  \begin{enumerate}
    \item Store a set of objects of known content and varied sizes, including
      bodies spanning many file-system blocks; leave one further store
      unfinished.
    \item Purge a subset of the stored objects, and apply metadata updates to
      another subset, both in place and withdrawing the record.
    \item Replay the log to each flush and FUA point in turn. For each replayed
      state: mount the file system (replaying its journal), construct the
      backend and let its restart procedure complete, look up every key, and
      read every hit's body in full.
    \item Repeat steps 1 to 3 with stores unsynchronised
      (\reqref{ADR-PERSISTENCE.UNSYNCED-01}), constructing the backend in
      each replayed state as after a host restart: under a boot identity
      other than the one the stores ran under.
  \end{enumerate}

  \textbf{Expected result:} In every replayed state: every hit's metadata and
  body are exactly one version the backend stored --- for an updated object,
  the version before or after an update, never a mixture; no object
  whose purge was acknowledged before that state's point is served; the
  unfinished store and every temporary are never served; every other object is
  either served as stored or a miss.

  \reqdiaglink{overview}{\textit{View the domain overview.}}

  \reqcover{satisfy}{IMP}

  \verify{HLR-PERSISTENCE.CRASHSAFE-01}
  \verify{FRE-COMPLIANCE.INCOMPLETE-01}
  \verify{FRE-COMPLIANCE.PURGEDURABLE-01}
  \verify{FRE-COMPLIANCE.METAPERSIST-01}
\end{requirement}

\begin{requirement}{TC-PERSISTENCE.CRASHDISCARD-01}[draft][verify]{Discard after a host restart}

  \localreqdiagram*{tccrashdiscard01}{\reqid}{\reqdesc}

  \textbf{Objective:} A cache that published stores without synchronising
  them serves none of its records after a host restart it was not flushed
  for, and all of them after a process restart or a host restart it was
  flushed for.

  \textbf{Setup:} The backend driven directly through the \texttt{Storage}
  interface on Linux, with a test hook that sets the boot identity the
  backend reads, and another that makes synchronising a store's temporary
  fail.

  \textbf{Steps:}
  \begin{enumerate}
    \item With stores unsynchronised and the synchronisation of a temporary
      made to fail, store objects, and leave one further store unfinished.
    \item Construct a second backend on the same cache root, with the same
      boot identity, and look up every key.
    \item Construct a third with a different boot identity, look up every key,
      and inspect the cache root.
    \item Store objects again, flush the backend, construct another with yet
      another boot identity, and look up every key.
    \item With stores strict, store objects, construct another backend with a
      different boot identity, and look up every key.
    \item Construct a backend with stores unsynchronised and the boot identity
      made unreadable.
  \end{enumerate}

  \textbf{Expected result:} In step 1 every finished store is published
  without reaching the failing synchronisation, and the instance's marker
  holds the boot identity. In step 2 every finished store is a hit. In step 3
  every key is a miss, no marker remains, and the discarded tree is gone once
  its background removal ends. In step 4 every object stored before the flush
  is a hit. In step 5 no marker is written and every object is a hit. In
  step 6 construction fails with an error naming the boot identity's source.

  \reqdiaglink{global}{\textit{View complete diagram.}}

  \reqcover{satisfy}{IMP}

  \verify{FDS-PERSISTENCE.CRASHDISCARD-01}
\end{requirement}

\begin{requirement}{TC-PERSISTENCE.RESTART-01}[draft][verify]{Objects survive a process exit}

  \localreqdiagram*{tc28}{\reqid}{\reqdesc}

  \textbf{Objective:} Every cached object whose storage completed before the
  cache process exited --- in an orderly way or by crashing --- is served by
  the process that starts next.

  \textbf{Setup:} The test-subject proxy with the file-system backend, its
  upstream pointed at an origin that counts the requests it receives.

  \textbf{Steps:}
  \begin{enumerate}
    \item Through the proxy, cache a set of objects: varied body sizes,
      \texttt{Vary} variants, and objects whose metadata was then updated by
      revalidation.
    \item Stop the proxy gracefully; start it again on the same cache
      directory; request every object.
    \item Cache a further set; as soon as the last store completes, kill the
      proxy; start it again; request every object from both sets.
  \end{enumerate}

  \textbf{Expected result:} After each restart every object is served from the
  cache, byte for byte and with its latest metadata, and the origin receives
  no request for any of them.

  \reqdiaglink{overview}{\textit{View the domain overview.}}

  \reqcover{satisfy}{IMP}

  \verify{HLR-PERSISTENCE.DURABLE-01}
\end{requirement}

\begin{requirement}{TC-PERSISTENCE.CAPACITY-01}[draft][verify]{Cache larger than host RAM}

  \localreqdiagram*{tc29}{\reqid}{\reqdesc}

  \textbf{Objective:} The cache holds, and serves, a set of objects whose total
  size exceeds the RAM installed on its host.

  \textbf{Setup:} The test-subject proxy with the file-system backend in a
  virtual machine with little RAM (e.g.\ 1~GiB) and ample disk, its eviction
  budget configured above the total the test stores.

  \textbf{Steps:}
  \begin{enumerate}
    \item Cache objects through the proxy until their total size is at least
      four times the machine's RAM, sampling the proxy's resident memory
      throughout.
    \item Restart the proxy and request every object.
  \end{enumerate}

  \textbf{Expected result:} Every store succeeds and no object is evicted;
  after the restart every object is served from the cache intact; the proxy is
  never killed for lack of memory, and its resident memory stays well below
  the machine's RAM.

  \reqdiaglink{overview}{\textit{View the domain overview.}}

  \reqcover{satisfy}{IMP}

  \verify{HLR-PERSISTENCE.CAPACITY-01}
\end{requirement}

% ---------------------------------------------------------------------
% Efficiency.
% ---------------------------------------------------------------------
\begin{requirement}{TC-EFFICIENCY.SPACE-01}[draft][verify]{Space overhead measured}

  \localreqdiagram*{tc33}{\reqid}{\reqdesc}

  \textbf{Objective:} Measure the cache's on-disk space overhead as
  \reqref{HLR-EFFICIENCY.SPACE-01} defines it and compare it with the 25\,\%
  target.

  \textbf{Setup:} The test-subject proxy with the file-system backend on a
  dedicated file system, its upstream an origin serving three stated loads:
  (a) a mix of object sizes recorded from, or modelled on, the intended
  deployment; (b) only objects below 4~KiB; (c) only objects of 1~MiB and
  more. A tool that walks the cache directory and sums, per record, the
  metadata segments and body lengths the record header gives.

  \textbf{Steps:}
  \begin{enumerate}
    \item For each load, fill the cache through the proxy until it holds at
      least $10^5$ objects, then let it settle (no store in flight).
    \item Measure the blocks the cache directory tree occupies (e.g.\
      \texttt{du} in bytes), and the objects' own metadata and body bytes with
      the tool; compute the overhead as their difference over the latter.
  \end{enumerate}

  \textbf{Expected result:} The overhead for each load is recorded. For the
  deployment load (a) it is at most 25\,\%. While
  \reqref{ADR-EFFICIENCY.COMPACTION-01}'s conflict stands, a result above the
  target for loads dominated by small objects is known non-conformance, to be
  reported with its measured value rather than treated as a regression.

  \reqdiaglink{overview}{\textit{View the domain overview.}}

  \reqcover{satisfy}{IMP}

  \verify{HLR-EFFICIENCY.SPACE-01}
\end{requirement}

\begin{requirement}{TC-PERFORMANCE.LOOKUP-01}[draft][verify]{Lookup latency flat with size}

  \localreqdiagram*{tcperflookup01}{\reqid}{\reqdesc}

  \textbf{Objective:} The 99th-percentile lookup latency with $10^7$ objects
  cached is no more than 25\,\% above that with $10^4$ objects cached, as
  \reqref{HLR-PERFORMANCE.LOOKUP-01} requires.

  \textbf{Setup:} A dedicated test host with the cache directory on an SSD and
  the file system a deployment would use, idle apart from the test. The
  backend is driven directly through the \texttt{Storage} interface, and the
  time from calling \texttt{lookup} to its return --- a hit with its metadata
  read, or a miss --- is recorded per call; the rest of the lookup path in
  \texttt{pingora-cache} does not depend on how many objects are stored. Objects carry
  uniformly distributed cache-key hashes and small bodies.

  \textbf{Steps:}
  \begin{enumerate}
    \item Store $10^4$ objects.
    \item \emph{Cold measurement:} drop the operating system's page cache,
      then perform at least $10^5$ lookups of uniformly chosen keys, half of
      them stored (hits) and half never stored (misses), recording each
      latency.
    \item \emph{Warm measurement:} repeat the same lookups immediately, so
      that the looked-up records and directories are in the page cache.
    \item Store further objects into the same cache directory until it holds
      $10^7$ objects, then repeat steps 2 and 3 with a fresh uniform key
      sample.
    \item For each of the four runs, compute the 99th-percentile latency
      separately for hits and for misses.
  \end{enumerate}

  \textbf{Expected result:} For hits and for misses, in both the cold and the
  warm condition, the 99th-percentile lookup latency at $10^7$ objects is no
  more than 1.25 times the corresponding value at $10^4$ objects.

  \reqdiaglink{overview}{\textit{View the domain overview.}}

  \reqcover{satisfy}{IMP}

  \verify{HLR-PERFORMANCE.LOOKUP-01}
\end{requirement}

\begin{requirement}{TC-PERFORMANCE.HITRATE-01}[draft][verify]{Hit throughput near in-memory}

  \localreqdiagram*{tcperfhitrate01}{\reqid}{\reqdesc}

  \textbf{Objective:} Cache hits on objects the host keeps in memory run at
  no less than 60\,\% of the request rate of the in-memory backend, for 1~KiB,
  64~KiB and 1~MiB bodies, as \reqref{HLR-PERFORMANCE.HITRATE-01} requires.

  \textbf{Setup:} A dedicated test host, idle apart from the test, with the
  cache directory on an SSD and the file system a deployment would use. One
  origin serves fixed bodies, marked cacheable. The fixture proxy runs once on
  this backend, in its default configuration, and once on \texttt{MemCache},
  each with one worker thread per processor it is given; both get the same
  processors, and the load generator gets other physical cores.
  \texttt{cargo xtask bench} is this harness.

  \textbf{Steps:}
  \begin{enumerate}
    \item For each body size --- 1~KiB, 64~KiB and 1~MiB --- and each
      backend, start the proxy with an empty cache, request each of 100
      objects of that size until every one is stored, then request them in
      random order over 64 concurrent connections for 10~s, recording the
      number of requests completed per second.
    \item Repeat for three rounds, changing the order of the backends each
      round, and take the median rate of each backend for each size.
  \end{enumerate}

  \textbf{Expected result:} For each body size, this backend's median rate is
  at least 0.60 times that of \texttt{MemCache}; no request fails, and the
  origin receives no request during the measured runs.

  \reqdiaglink{global}{\textit{View complete diagram.}}

  \reqcover{satisfy}{IMP}

  \verify{HLR-PERFORMANCE.HITRATE-01}
\end{requirement}

\begin{requirement}{TC-PERFORMANCE.STORERATE-01}[draft][verify]{Store rate with host-crash loss}

  \localreqdiagram*{tcperfstorerate01}{\reqid}{\reqdesc}

  \textbf{Objective:} With stores unsynchronised, storing new objects runs at
  no less than 40\,\% of the request rate of the in-memory backend, for 4~KiB
  and 64~KiB bodies, as \reqref{HLR-PERFORMANCE.STORERATE-01} requires.

  \textbf{Setup:} As for \reqref{TC-PERFORMANCE.HITRATE-01}, with this backend
  configured for unsynchronised stores (\reqref{ADR-PERSISTENCE.UNSYNCED-01}).

  \textbf{Steps:}
  \begin{enumerate}
    \item For each body size --- 4~KiB and 64~KiB --- and each backend, start
      the proxy with an empty cache and request objects never requested
      before over 64 concurrent connections, 20{,}000 of 4~KiB or 10{,}000 of
      64~KiB, recording the number of requests completed per second.
    \item Repeat for three rounds, changing the order of the backends each
      round, and take the median rate of each backend for each size.
  \end{enumerate}

  \textbf{Expected result:} For each body size, this backend's median rate is
  at least 0.40 times that of \texttt{MemCache}; no request fails, and the
  origin receives every request once.

  \reqdiaglink{global}{\textit{View complete diagram.}}

  \reqcover{satisfy}{IMP}

  \verify{HLR-PERFORMANCE.STORERATE-01}
\end{requirement}

% ---------------------------------------------------------------------
% Pingora compatibility.
% ---------------------------------------------------------------------
\begin{requirement}{TC-COMPAT.EVICTION-01}[draft][verify]{Exact eviction state after restart}

  \localreqdiagram*{tc11}{\reqid}{\reqdesc}

  \textbf{Objective:} Once the backend reports its restart procedure
  complete, the eviction manager's view of admitted objects is consistent
  with the durably published objects.

  \textbf{Setup:} The backend with the stock \texttt{simple\_lru} eviction
  manager and Eviction Recovery, driven by a test harness that can read the
  manager's entries and totals.

  \textbf{Steps:}
  \begin{enumerate}
    \item Populate the cache with a known set of objects, including
      \texttt{Vary} variants and objects of very different sizes.
    \item Before restarting, also leave on disk: an object purged during the
      previous run, a record damaged out of band, and a crash-orphaned
      per-writer temporary; and remove one variant's primary.
    \item Restart the process and wait until the backend reports its restart
      procedure complete.
    \item Inspect the manager's \texttt{total\_items} and \texttt{total\_size} and
      its per-entry view; compare each entry's size with the length of the
      file that stores the object.
    \item Store one more object; confirm the size it is admitted with is also
      its record file length.
  \end{enumerate}

  \textbf{Expected result:} Every durably published object, the variant whose
  primary was removed included, is represented exactly once with its record
  size and freshness deadline; no entry exists for the purged object, the
  damaged record, the temporary, or the removed primary; the reported totals
  equal the count and aggregate record size of the
  published set; a store completed after the restart is admitted with the same
  size definition.

  \reqdiaglink{overview}{\textit{View the domain overview.}}

  \reqcover{satisfy}{IMP}

  \verify{FRE-COMPAT.EVICTION-01}
\end{requirement}

\begin{requirement}{TC-COMPAT.HOTUPGRADE-01}[draft][verify]{Cache consistent across hot upgrade}

  \localreqdiagram*{tc12}{\reqid}{\reqdesc}

  \textbf{Objective:} Once the outgoing process receives the Pingora upgrade
  signal, its backend makes no further durable change and refuses purges and
  metadata updates; the incoming process's activity never makes unreadable an
  object the outgoing process holds open; and the hot-upgrade helper starts
  the incoming process's restart scan only once the server has taken over
  the listening sockets and the close timeout has elapsed.

  \textbf{Setup:} Two builds of the test-subject proxy with the file-system
  backend and its hot-upgrade helper installed, run as real Pingora server
  processes sharing one cache directory. From step 6 on, a Pingora server
  run inside the test harness with the helper installed, the harness
  standing in for the outgoing process.

  \textbf{Steps:}
  \begin{enumerate}
    \item Run an old process serving normally, with streaming responses in
      progress against existing objects and a store of a slow upstream
      response in flight.
    \item Start a new process and send the old one the Pingora upgrade signal
      (\texttt{SIGQUIT}), letting the listening sockets transfer.
    \item Through the old process, while it still accepts or drains: let the
      in-flight store complete, and request a purge and trigger a
      revalidation metadata update for existing keys.
    \item Through the new process: purge an object the old process is
      streaming, and drive stores, metadata updates, and eviction.
    \item Let the old process finish draining and exit; restart nothing, and
      inspect the cache directory and the new process's lookups.
    \item Start a server that takes over its listening sockets from an
      outgoing process, bootstrapping as a service, with the helper given
      that service's handle; hand a listening socket over only some time
      after the server has started. Then start an ordinary server, taking
      over no socket, with the helper installed.
    \item Send one server with the helper installed the upgrade signal,
      though it has no socket to hand over, and another a graceful
      terminate.
    \item Start the helper's recovery service, and shut its server down
      while the service waits to scan; then, over other caches, once the
      scan has admitted its first record, and once its walk is done, with
      a record removed meanwhile.
  \end{enumerate}

  \textbf{Expected result:} The store the old process completed after the
  signal is not durably published; its purge and metadata update fail with an
  error and change nothing on disk; every response it was streaming completes
  byte for byte, including the one whose object the new process purged; the
  new process's lookups see only objects published before the signal or by
  the new process; the restart scan of step 6's first server begins no
  sooner than the close timeout after the hand-over, and the ordinary
  server's without waiting for it; in step 7, the upgrade signal makes
  its server's backend read-only, and the terminate leaves the other's
  writable; in step 8, the first service starts no scan and says so, the
  second scan admits no further record, and the third skips its ghost-entry
  checks, keeping the removed record's entry, both summary lines saying the
  server was shutting down.

  \reqdiaglink{overview}{\textit{View the domain overview.}}

  \reqcover{satisfy}{IMP}

  \verify{FRE-COMPAT.HOTUPGRADE-01}
  \verify{FDS-COMPAT.READONLY-01}
\end{requirement}

\begin{requirement}{TC-COMPAT.STREAMING-01}[draft][verify]{Followers stream a partial hit}

  \localreqdiagram*{tc16}{\reqid}{\reqdesc}

  \textbf{Objective:} The backend advertises streaming partial writes
  (\reqref{FRE-COMPAT.STREAMFLAG-01}), and while a leader writes an object,
  concurrent requests are served by partial hits bound to the most recently
  opened writer when several are in flight (\reqref{FRE-COMPAT.STREAMING-01}),
  streamed within the chunk-delivery latency bound
  (\reqref{FRE-COMPAT.STREAMDELIVERY-01}) and ending normally when the store
  completes (\reqref{FRE-COMPAT.STREAMEND-01}), from a single upstream fetch.

  \textbf{Setup:} A proxy with the file-system backend and no
  \texttt{CacheLock} configured, so that a second same-key store can open
  while the first is still in flight (with a lock configured, the same
  overlap needs the lock's age timeout to expire), on a test host that is not
  CPU-saturated, with at least ten followers and a body of at least a
  thousand chunks, so the delivery percentile rests on enough samples.

  \textbf{Steps:}
  \begin{enumerate}
    \item Query the backend's \texttt{support\_streaming\_partial\_write}.
    \item On a second proxy instance identical but for a \texttt{CacheLock}
      being configured, drive a cache miss from a slow upstream; once the
      miss handler opens, and before the body is written, record whether
      the cache lock for the key is still held.
    \item Drive a cache miss from an upstream that emits a large body in
      throttled chunks; let the miss handler open.
    \item While that store is still in flight, start a second store under
      the same key from a distinct upstream response, staggered so both
      writers' write windows overlap; let its miss handler open too.
    \item Issue a lookup for the key now that both writers are open; confirm
      the returned partial hit is backed by the second writer's --- the most
      recently opened --- temporary object, not the first's.
    \item Issue further concurrent lookups for the same key while the second
      store is in flight; confirm each returns a partial hit backed by that
      same temporary object.
    \item Observe the followers' streams as the second writer appends,
      timestamping each chunk's \texttt{write\_body} return and each
      follower's \texttt{read\_body} return with that chunk, before that
      store completes.
    \item Let the second store complete; drain the followers' streams
      (across the concluding atomic rename) and compare the received bodies
      to its upstream body.
    \item Count upstream fetches: exactly one for each writer, including the
      first, whose output was superseded before ever being streamed to a
      follower.
  \end{enumerate}

  \textbf{Expected result:} \texttt{support\_streaming\_partial\_write}
  returns \texttt{true}, and the cache lock is released once the leader's
  miss handler opens, before its body is written. A lookup issued while
  several same-key writes are in flight resolves to the most recently opened one, never an earlier
  writer. Followers receive body chunks as that writer writes them, with at
  least 99\,\% of all chunk deliveries, across every follower, within 10~ms
  of the chunk's \texttt{write\_body} return --- no follower waits for the
  whole object; every follower's body is
  byte-identical to its upstream response and ends with a normal end of
  body, including for followers still reading when the object is renamed to
  its canonical name; each writer's upstream is fetched exactly
  once.

  \reqdiaglink{overview}{\textit{View the domain overview.}}

  \reqcover{satisfy}{IMP}

  \verify{FRE-COMPAT.STREAMFLAG-01}
  \verify{FRE-COMPAT.STREAMING-01}
  \verify{FRE-COMPAT.STREAMDELIVERY-01}
  \verify{FRE-COMPAT.STREAMEND-01}
  \verify{FDS-STORAGE.PARTIAL-01}
\end{requirement}

\begin{requirement}{TC-COMPAT.UNSIZED-01}[draft][verify]{Unsized body accepted and persisted}

  \localreqdiagram*{tc19}{\reqid}{\reqdesc}

  \textbf{Objective:} A body whose total length is not declared when the
  store begins is persisted completely, and every chunk is in the store's
  temporary before its \texttt{write\_body} call returns, so the backend
  holds no body bytes beyond the chunk being written, as
  \reqref{FRE-COMPAT.UNSIZED-01} requires.

  \textbf{Setup:} The backend driven directly through the \texttt{Storage}
  interface by a test harness that can read the store's temporary in the
  cache directory.

  \textbf{Steps:}
  \begin{enumerate}
    \item Open a store and stream a body in irregular, unevenly-sized chunks
      without declaring its length, so the total length is unknowable when
      the store opens.
    \item After each \texttt{write\_body} call returns, compare the
      temporary's length with the length of its header, key, and metadata
      region plus the body bytes streamed so far.
    \item Once the stream ends, retrieve the object and compare its persisted
      length and bytes to the total streamed.
    \item Repeat with a body whose size happens to be zero (empty), and with
      bodies large enough to span multiple write chunks and file-system block
      boundaries.
  \end{enumerate}

  \textbf{Expected result:} Every store succeeds; after every
  \texttt{write\_body} call the temporary holds exactly the header, key,
  metadata region, and the body bytes streamed so far; the persisted
  object's length and bytes equal exactly what was streamed in every case,
  including the empty-body and multi-chunk boundary cases.

  \reqdiaglink{overview}{\textit{View the domain overview.}}

  \reqcover{satisfy}{IMP}

  \verify{FRE-COMPAT.UNSIZED-01}
\end{requirement}

\begin{requirement}{TC-COMPAT.CONCURRENT-01}[draft][verify]{Same-key writes stay isolated}

  \localreqdiagram*{tc20}{\reqid}{\reqdesc}

  \textbf{Objective:} Concurrent writes of the same key never interleave
  their bytes, and each is durably published as a complete object or not
  at all, while in-process lookups may see a partial hit of one in-flight
  writer. (Convergence on a single winner is covered separately by
  \reqref{TC-COMPLIANCE.REPLACE-01}.)

  \textbf{Setup:} The backend driven directly through the \texttt{Storage}
  interface by a test harness; a ``restart'' constructs a new backend instance
  over the same cache directory, with a second backend instance in a separate
  process for the cross-process case.

  \textbf{Steps:}
  \begin{enumerate}
    \item Start two concurrent stores under the same key, each with a
      distinctive, easily distinguishable byte pattern and a different
      length, staggered so their write windows overlap.
    \item While both stores are in flight, issue concurrent lookups for the
      key.
    \item Let both stores run to completion.
    \item Repeat with same-key writers started within the same process (not
      only across processes), exercising the same-process race the cache
      lock's age timeout permits.
  \end{enumerate}

  \textbf{Expected result:} No lookup ever observes a body mixing bytes from
  more than one writer; every lookup either returns one writer's object whole
  (or a partial hit of one in-flight writer, per
  \reqref{FRE-COMPAT.STREAMING-01}) or reports a miss --- never a torn
  or interleaved result. Once both stores complete, the backend has converged
  on exactly one whole, uncorrupted object.

  \reqdiaglink{overview}{\textit{View the domain overview.}}

  \reqcover{satisfy}{IMP}

  \verify{FRE-COMPAT.CONCURRENT-01}
\end{requirement}

\begin{requirement}{TC-COMPAT.FORMATUPGRADE-01}[draft][verify]{Foreign records are misses}

  \localreqdiagram*{tc30}{\reqid}{\reqdesc}

  \textbf{Objective:} A record written under another record format version is
  reported as a miss and never decoded, while a record written under the
  current one is returned exactly; and the \texttt{pingora-cache} release in
  use encodes metadata exactly as the current format version relies on.

  \textbf{Setup:} The backend driven directly through the \texttt{Storage}
  interface, instrumentation recording every metadata decode, and a fixed
  \texttt{CacheMeta} together with the bytes \texttt{pingora-cache} 0.9.0
  encodes it to.

  \textbf{Steps:}
  \begin{enumerate}
    \item Construct the backend; store a set of objects.
    \item Construct it again over the same directory; look every object up.
    \item Store a fresh set; rewrite each record's \texttt{version} field to
      another value, resealing its \texttt{header\_crc}; construct the backend
      again; look every object up; inspect the cache directory once its
      restart procedure completes.
    \item Encode the fixed \texttt{CacheMeta} with the \texttt{pingora-cache}
      in use and compare the result with the fixed bytes; decode the fixed
      bytes and compare every field with the fixed \texttt{CacheMeta}.
  \end{enumerate}

  \textbf{Expected result:} Step 2 returns every object exactly as stored.
  Step 3 reports every rewritten object as a miss, and the instrumentation
  shows no attempt to decode their metadata; the foreign records are
  reclaimed by the restart procedure (\reqref{FDS-STORAGE.RECLAIM-01}). In
  step 4 both comparisons are equal. A difference means the release in use
  changed the encoding the current format relies on, and the format version
  must be reconsidered before that release is adopted.

  \reqdiaglink{overview}{\textit{View the domain overview.}}

  \reqcover{satisfy}{IMP}

  \verify{FRE-COMPAT.FORMATUPGRADE-01}
\end{requirement}

\begin{requirement}{TC-COMPAT.WRITEFAIL-01}[draft][verify]{Write failures fail the store}

  \localreqdiagram*{tc31}{\reqid}{\reqdesc}

  \textbf{Objective:} A failure writing or synchronising a store's data is
  returned as an error from the call during which it occurred, and nothing is
  published.

  \textbf{Setup:} The backend driven directly through the \texttt{Storage}
  interface, with its cache directory on a small file system (e.g.\ a
  loop-mounted image of a few megabytes) over a device-mapper
  \texttt{flakey} target able to fail writes on demand.

  \textbf{Steps:}
  \begin{enumerate}
    \item Stream a body larger than the file system's free space.
    \item On a file system with space, stream a complete body, switch the
      device to failing writes, and call \texttt{finish}.
    \item Remove, from outside the backend, a shard directory the backend has
      already created; store an object into that shard.
    \item Remove that shard directory again and make its first-level parent
      directory read-only, so it cannot be recreated; store another object
      into it.
    \item After each case, look the key up.
  \end{enumerate}

  \textbf{Expected result:} Step 1's \texttt{write\_body} and step 2's
  \texttt{finish} each return an error; step 3's store succeeds after the
  backend recreates the directory once; step 4's store returns an error; no
  failed store is ever published, so each such lookup reports a miss.

  \reqdiaglink{overview}{\textit{View the domain overview.}}

  \reqcover{satisfy}{IMP}

  \verify{FRE-COMPAT.WRITEFAIL-01}
  \verify{FDS-LAYOUT.STARTUP-01}
\end{requirement}

\begin{requirement}{TC-COMPAT.ABORTSPACE-01}[draft][verify]{Aborted stores release space}

  \localreqdiagram*{tc32}{\reqid}{\reqdesc}

  \textbf{Objective:} A store that ends without being published loses its
  temporary no later than when its miss handler is dropped, and the storage
  it consumed is freed without a restart, however it ended.

  \textbf{Setup:} The backend driven directly through the \texttt{Storage}
  interface, as in \reqref{TC-COMPAT.WRITEFAIL-01}.

  \textbf{Steps:}
  \begin{enumerate}
    \item Stream part of a large body, then drop the miss handler without
      finishing, as \texttt{HttpCache} does after an upstream failure or a
      cancellation; inspect the cache directory and free space.
    \item Repeat with a partial-hit reader attached to the store before the
      handler is dropped; read its stream to the end.
    \item Repeat with a store ended by a write failure.
    \item Repeat step 1 a thousand times in a row.
  \end{enumerate}

  \textbf{Expected result:} After each dropped handler the store's temporary is
  gone at once, and the free space it held is available again once the
  backend has closed the file, without a restart; the attached reader's
  stream ends with an error; after the thousand aborts the cache directory's
  disk usage is back to what it was before them.

  \reqdiaglink{overview}{\textit{View the domain overview.}}

  \reqcover{satisfy}{IMP}

  \verify{FRE-COMPAT.ABORTSPACE-01}
\end{requirement}

% ---------------------------------------------------------------------
% Eviction-state reconstruction mechanism.
% ---------------------------------------------------------------------
\begin{requirement}{TC-EVICTION.RECONSTRUCT-01}[draft][verify]{Reconstruction converges exactly}

  \localreqdiagram*{tc13}{\reqid}{\reqdesc}

  \textbf{Objective:} When \texttt{recover()} returns, the manager reflects
  the published set exactly, with record sizes --- whether serving waited for
  it or ran alongside it --- with no snapshot or delta log involved.

  \textbf{Setup:} The backend with the stock \texttt{simple\_lru} eviction
  manager and Eviction Recovery behind an \texttt{HttpCache}, driven by a test
  harness with hooks that pause \texttt{recover()} at chosen points.

  \textbf{Steps:}
  \begin{enumerate}
    \item Populate a large known set of objects and restart the process;
      call \texttt{recover()} without serving, and let it return.
    \item Restart again, and this time begin serving at once. With test
      hooks pausing \texttt{recover()} at chosen points, drive live traffic
      that: replaces an object after the scan has read it but before it is
      admitted; removes an object's record after the scan admitted it,
      leaving its entry in the manager, as a live eviction's purge does when
      it lands after the scan admitted its victim again; republishes a key
      after the scan admitted it and a purge removed it;
      and stores objects the scan never reaches.
    \item Let \texttt{recover()} return; compare the manager's entries,
      sizes, \texttt{total\_items} and \texttt{total\_size} with the files on
      disk.
    \item Trace file-system activity throughout both runs.
    \item Make the backend read-only, then call \texttt{recover()}, with a
      test hook pausing it after it reads any record.
  \end{enumerate}

  \textbf{Expected result:} After each run, every published object is
  represented exactly once with its current record size, no entry exists for
  a removed object, and the totals equal the count and aggregate record size
  of the published set --- including after each interleaving in step 2; no
  snapshot file or append-only log is read or written at any point; a scan of
  a read-only backend returns without reading any record or admitting
  anything.

  \reqdiaglink{overview}{\textit{View the domain overview.}}

  \reqcover{satisfy}{IMP}

  \verify{FDS-EVICTION.STARTUP-01}
\end{requirement}

\begin{requirement}{TC-STORAGE.PUBLISH-01}[draft][verify]{Atomic same-directory publication}

  \localreqdiagram*{tc21}{\reqid}{\reqdesc}

  \textbf{Objective:} The per-writer temporary and rename publish atomically,
  keep concurrent writers isolated, never reuse another writer's file,
  synchronise before publishing, report the record size, and publish nothing
  once the backend is read-only.

  \textbf{Setup:} The backend driven directly through the \texttt{Storage}
  interface, with system calls traced (e.g.\ with \texttt{strace}).

  \textbf{Steps:}
  \begin{enumerate}
    \item Start two concurrent same-key writers; confirm each writes its own
      temporary, named with the writer identity and counter
      \reqref{FDS-STORAGE.PUBLISH-01} specifies, and neither's bytes appear in
      the other's file.
    \item Before a store opens, create a file at exactly the temporary name
      its counter will produce next; open the store.
    \item Let a store run to completion; confirm the object becomes visible
      under its canonical name only at the rename, and inspect the traced
      system-call order and the size \texttt{finish} reports.
    \item Crash a writer mid-store, before rename; confirm the abandoned
      temporary is never served or promoted under the canonical name; restart
      the process and confirm the temporary is gone after its restart-time
      scan.
    \item Trigger a replacement store under an existing key; confirm a
      reader already streaming the outgoing object completes undisturbed
      while the new object publishes, and subsequent lookups return only
      the new object.
    \item Make the backend read-only (\reqref{FDS-COMPAT.READONLY-01}) while a
      store is in flight; let it complete.
    \item Inspect the implementation and the tree layout: every temporary is
      created in the directory of its canonical name, and no code path
      publishes by copying.
  \end{enumerate}

  \textbf{Expected result:} Concurrent writers' bytes are never interleaved;
  the pre-existing file at the temporary name is left untouched and the store
  uses the next counter value; the temporary's data is synchronised before the
  rename, and \texttt{finish} reports the published file's full length; a
  crash leaves no promoted incomplete object; the rename is the sole point of
  visibility; replacement converges to a single winner without disturbing
  in-flight readers of the old object; the store completing once the backend
  is read-only
  reports success, publishes nothing, and leaves no temporary; the inspection
  finds no cross-directory staging and no copy-based publication.

  \reqdiaglink{overview}{\textit{View the domain overview.}}

  \reqcover{satisfy}{IMP}

  \verify{FDS-STORAGE.PUBLISH-01}
\end{requirement}

% ---------------------------------------------------------------------
% Purge: canonical-name unlink only.
% ---------------------------------------------------------------------
\begin{requirement}{TC-STORAGE.PURGE-01}[draft][verify]{Purge unlinks canonical name only}

  \localreqdiagram*{tc25}{\reqid}{\reqdesc}

  \textbf{Objective:} Purge resolves and unlinks only the canonical path
  derived from the presented target key's combined hash, reports the outcome
  the framework distinguishes, is idempotent and a harmless no-op
  on an absent key, ends a same-key store in flight when it invalidates and
  leaves it alone when it evicts, and disturbs no reader already streaming
  the file.

  \textbf{Setup:} The backend driven directly through the \texttt{Storage}
  interface by a test harness; a ``restart'' constructs a new backend instance
  over the same cache directory.

  \textbf{Steps:}
  \begin{enumerate}
    \item Store an object together with a \texttt{Vary} variant under the
      same primary key; purge the primary; confirm only the primary's
      canonical file is removed from disk and the variant's own file is
      untouched.
    \item Store a primary and a variant again; purge using the variant's
      own (varianced) \texttt{CompactCacheKey}, as
      \texttt{HttpCache::spawn\_async\_purge} can present it; confirm only
      the variant's canonical file is removed and the primary's is
      untouched.
    \item Purge the same key a second time, and purge a key that was never
      stored; confirm both report \texttt{NotFound}, with no error and no
      side effect on any other object.
    \item Purge a stored key through each target kind in turn --- an
      \texttt{Active} target, and an \texttt{Exact} target holding a
      \texttt{KeyOnly} entry key --- and confirm each removes that key's
      canonical file and reports \texttt{Purged(None)}, never an entry
      identifier (\reqref{ADR-COMPAT.ENTRYID-01}).
    \item Purge a stored key through an \texttt{Exact} target holding an
      \texttt{Identified} entry key; confirm it reports \texttt{NotFound}
      and the file is still on disk.
    \item Start a same-key store and let it reach the point of holding an
      open per-writer temporary, and attach a partial-hit reader to it;
      invalidate the key; look the key up, with and without the store's
      streaming write tag; let the store complete; look the key up; then
      store the key afresh. Invalidate a key that holds nothing but a store
      in flight. Repeat the first case with a purge made to evict instead of
      the invalidation.
    \item Attach a reader mid-stream on an already-published object; purge
      its key while the read is in progress; confirm the reader runs to
      the end of the body it already opened, never an error and never
      another object's bytes, and a subsequent lookup reports a miss.
    \item Race an invalidation, and then a purge made to evict, against a
      same-key store's completion, from both orderings: the purge before the
      store's rename, and after it.
    \item With system calls traced, purge an existing object; then, with
      directory synchronisation made to fail (e.g.\ over a device-mapper
      \texttt{flakey} target), purge another. With synchronisation still
      failing, evict an existing object, then invalidate the evicted key,
      and invalidate a key whose shard directory does not exist.
    \item Make the backend read-only (\reqref{FDS-COMPAT.READONLY-01}) and
      purge an existing object.
    \item Put a directory where a key's record would be, so that the
      \texttt{unlink} fails whatever the test's privileges, and purge the
      key.
  \end{enumerate}

  \textbf{Expected result:} Only the presented target key's canonical path is
  ever named or opened by purge; a primary purge never touches a \texttt{Vary}
  variant, a variant purge never touches its primary, and a purge never
  touches a writer's own temporary; an \texttt{Active} and a
  \texttt{KeyOnly} \texttt{Exact} target resolve identically and report
  \texttt{Purged(None)}, while an \texttt{Identified} \texttt{Exact} target
  reports \texttt{NotFound} without touching the file; repeated and
  never-stored purges report \texttt{NotFound} harmlessly; a concurrent reader completes or fails cleanly, never
  crossing into another object's bytes; an invalidation ends a same-key
  store in flight: an untagged lookup then misses while a tagged one still
  finds the store, the reader attached to it gets the whole body, and the
  store's completion publishes nothing and reports an error, its temporary
  removed, while a store begun after the purge publishes as usual; the
  invalidation of a key holding only a store in flight reports
  \texttt{Purged(None)}; a purge made to evict leaves the store to publish;
  racing a store, an invalidation leaves the key absent in either order,
  and a purge made to evict leaves whichever of the two landed last; an
  invalidation
  synchronises the shard directory after the unlink, or after finding
  nothing, before it reports, and fails with an error when the
  synchronisation does, while a purge made to evict does not synchronise and
  so still removes its record, and an invalidation whose shard directory does
  not exist has nothing to synchronise and reports \texttt{NotFound};
  the purge once the backend is read-only fails with an error and leaves
  the file in place; and the purge whose \texttt{unlink} fails reports an
  error and leaves the path as it was.

  \reqdiaglink{overview}{\textit{View the domain overview.}}

  \reqcover{satisfy}{IMP}

  \verify{FDS-STORAGE.PURGE-01}
\end{requirement}

% ---------------------------------------------------------------------
% Reclaim: crash-orphaned temporaries, and damaged and withdrawn records.
% ---------------------------------------------------------------------
\begin{requirement}{TC-STORAGE.EXPIRE-01}[draft][verify]{Expire keeps the body readable}

  \localreqdiagram*{tc34}{\reqid}{\reqdesc}

  \textbf{Objective:} An expiry stops the object being served as a fresh hit
  while leaving its body byte-identical and readable for the revalidation that
  follows, takes the in-place metadata path whatever headroom the record has
  left, and reports the outcome each case calls for.

  \textbf{Setup:} The backend driven directly through the \texttt{Storage}
  interface by a test harness, with the record's file checked so a write in
  place can be told from a new file, and with the metadata headroom
  configurable per instance (\reqref{SDS-LAYOUT.OBJECTFILE-01}).

  \textbf{Steps:}
  \begin{enumerate}
    \item Store an object with a known body and a freshness lifetime well in
      the future; expire its key; look the key up and read the body to the
      end.
    \item Expire the same key a second time, and expire a key that was never
      stored.
    \item Expire a stored key through an \texttt{Exact} target holding an
      \texttt{Identified} entry key.
    \item With headroom configured, expire an object and compare the file's
      identity (device and inode) and length before and after; then repeat
      against an instance configured with zero headroom, for an object whose
      freshness deadline falls on a whole second, expired at an instant well
      inside a second.
    \item Store an object whose response sets a
      \texttt{stale-while-revalidate} window; expire it; inspect the metadata
      the next lookup returns.
    \item Start a same-key store and let it reach the point of holding an open
      per-writer temporary; expire the key concurrently; let the store
      complete; look the key up.
    \item Attach a reader mid-stream on an already-published object; expire its
      key while the read is in progress.
    \item Pause an expiry between reading the record's metadata and writing it
      back; from another task, run a metadata update on the same key; resume the
      expiry; inspect the resulting metadata.
    \item Make the backend read-only (\reqref{FDS-COMPAT.READONLY-01}) and
      expire an existing object.
    \item Expire a stored key whose record has been damaged, truncated,
      rewritten under another format version, or withdrawn
      (\reqref{SDS-LAYOUT.OBJECTFILE-01}), comparing the record's bytes
      before and after; look the key up.
  \end{enumerate}

  \textbf{Expected result:} After an expiry the lookup still finds the record
  and its body reads back byte-identical, while the metadata it returns is no
  longer fresh at any instant after the expiry; a second expiry is
  idempotent and a never-stored key reports \texttt{NotFound}, neither an error
  nor a side effect; an \texttt{Identified} \texttt{Exact} target reports
  \texttt{NotFound} and leaves the record untouched; the expiry writes the
  metadata in place, leaving the file's device, inode and length unchanged,
  with headroom and without, the whole-second deadline included; the
  object's serve-stale window is still open after the expiry, so the stored
  body remains servable while the revalidation runs; the concurrent store's
  completing rename replaces the expired record with its fresh object; a
  reader streaming the object completes its body undisturbed; the paused
  expiry and the concurrent metadata update serialise on the publication
  lock, so the record ends carrying one of the two complete metadata sets
  and never a mixture;
  the expiry once the backend is read-only fails with an error, leaving the
  stored metadata unchanged; and each unreadable record is reported
  \texttt{NotFound} and left byte for byte as it was, and a lookup of it
  misses.

  \reqdiaglink{overview}{\textit{View the domain overview.}}

  \reqcover{satisfy}{IMP}

  \verify{FRE-COMPAT.EXPIRE-01}
  \verify{FDS-STORAGE.EXPIRE-01}
\end{requirement}

\begin{requirement}{TC-STORAGE.EXECUTOR-01}[draft][verify]{Storage never blocks the worker}

  \localreqdiagram*{tc35}{\reqid}{\reqdesc}

  \textbf{Objective:} No \texttt{Storage} call, hit-handler read, or
  miss-handler write holds the executor thread that polls it for the
  operation's duration, nor does the restart scan, so a request sharing that
  worker is not delayed by another request's file I/O.

  \textbf{Setup:} The backend driven directly through the \texttt{Storage}
  interface on a \emph{single-worker} runtime --- the Pingora default for a
  service --- with a second task on the same worker ticking on a short timer
  and recording the longest gap between its ticks. That gap is what a
  co-resident request would have waited.

  \textbf{Steps:}
  \begin{enumerate}
    \item Start the ticking task and let it settle.
    \item Drive a sequence of complete cache operations against bodies large
      enough that reads and writes are real work: store, look up and read the
      body to its end, update the metadata, expire, and purge.
    \item Open a store and attach a partial-hit reader to it, then alternate
      \texttt{write\_body} and \texttt{read\_body} so the write and read
      paths are both active, and finish the store.
    \item Read the longest inter-tick gap after each scenario.
    \item Repeat both scenarios against a build whose blocking work runs
      inline on the worker instead of on the blocking pool.
    \item Run Eviction Recovery (\reqref{FDS-EVICTION.STARTUP-01}) on a
      single worker over a damaged record, a stopped writer's temporary, and
      a record the eviction manager evicts as soon as it is admitted. With a
      test hook holding each removal where it frees its file, ask another
      task on that worker to answer.
    \item Run Eviction Recovery on a single worker beside another task that
      notes each turn it gets, over records whose files are all removed
      between the walk and the ghost-entry check, and note each call the
      scan makes into the eviction manager.
  \end{enumerate}

  \textbf{Expected result:} In both scenarios the longest gap stays within a
  bound far below what one operation's file I/O costs, showing the worker was
  free throughout; the inline build exceeds it in both, so the measurement is
  shown to detect the condition it exists to rule out. While the restart
  scan frees a file it removed, the other task on its worker answers, and it
  gets a turn between any two calls the scan makes into the eviction manager
  --- one admission and one removal per record. The bound is a
  \emph{stall} bound, not a latency figure for the operations themselves,
  which \reqref{TC-PERFORMANCE.LOOKUP-01} measures.

  \reqdiaglink{overview}{\textit{View the domain overview.}}

  \reqcover{satisfy}{IMP}

  \verify{FDS-STORAGE.EXECUTOR-01}
\end{requirement}

\begin{requirement}{TC-STORAGE.PAGECACHE-01}[draft][verify]{Page-cache reads on the worker}

  \localreqdiagram*{tcpagecache01}{\reqid}{\reqdesc}

  \textbf{Objective:} On Linux, the lookup and the body reads of a published
  record the page cache holds complete on the worker, without a handoff to
  the blocking pool; each falls back to the pool when a call would wait, and
  the answer is the same either way.

  \textbf{Setup:} The backend driven directly through the \texttt{Storage}
  interface on Linux, with test hooks that count the operations handed to the
  blocking pool, and that make the worker's open, or a given read, report
  that it would wait.

  \textbf{Steps:}
  \begin{enumerate}
    \item Store records with bodies of 1~KiB and of 200~KiB, the second
      spanning several chunks; look each up and read its body to the end,
      counting handoffs.
    \item Make the worker's open report a wait, and repeat step 1.
    \item Make the second body read of the larger record report a wait, and
      repeat step 1.
    \item Look up a key never stored, and a record whose metadata checksum
      does not match.
    \item Switch page-cache reads off in the configuration, and repeat
      step 1.
  \end{enumerate}

  \textbf{Expected result:} In step 1 no operation is handed to the pool. In
  step 2 each lookup is handed over once, and in step 3 the larger record's
  second chunk once. In step 4 both lookups are misses, the damaged record's
  reached through the pool's locked re-read. In step 5 every lookup and every
  chunk is handed over once, as \reqref{FDS-STORAGE.EXECUTOR-01} specifies.
  In every step, each hit's metadata and body are exactly those stored.

  \reqdiaglink{global}{\textit{View complete diagram.}}

  \reqcover{satisfy}{IMP}

  \verify{FDS-STORAGE.PAGECACHE-01}
\end{requirement}

\begin{requirement}{TC-STORAGE.RECLAIM-01}[draft][verify]{Orphan reclaim at restart}

  \localreqdiagram*{tc26}{\reqid}{\reqdesc}

  \textbf{Objective:} The restart-time scan reclaims a temporary exactly when
  its writer is not running, with that writer's lease, and a damaged, foreign
  or withdrawn record --- never admitting what it reclaims, never reclaiming a
  \texttt{Vary} variant, whatever its primary holds, never unlinking a record
  a live store has just replaced, and reporting what it reclaimed.

  \textbf{Setup:} The backend with the stock \texttt{simple\_lru} eviction
  manager and Eviction Recovery, driven by a test harness with a hook that
  pauses the scan between classifying and unlinking an entry.

  \textbf{Steps:}
  \begin{enumerate}
    \item Leave temporaries named after writers that are not running: one
      whose lease does not exist, and one whose lease exists but is not
      locked, as a crashed writer leaves it. Construct the backend (running
      its restart procedure); confirm both temporaries are unlinked, and so
      is the unlocked lease.
    \item Leave a temporary named after a lease the harness holds locked, as
      a live peer would; construct the backend; confirm the temporary and
      its lease are left untouched, then confirm a later scan reclaims both
      once the harness releases the lease.
    \item Leave a record damaged out of band (as in
      \reqref{TC-LAYOUT.OBJECTFILE-01}), a record of another format version,
      and a withdrawn record (\reqref{FDS-STORAGE.METAUPDATE-01}); also copy
      one key's record to the path of a key never stored, and another's,
      under its own name, into a shard directory that name does not select;
      construct the backend; confirm all five are unlinked, the withdrawn one
      counted under its own class and the others as damaged, none is
      admitted, and the records at their own paths are.
    \item Store four primaries, each with a \texttt{Vary} variant. Remove the
      first primary's file out of band (simulating a completed purge or
      eviction), damage the second's (as in
      \reqref{TC-LAYOUT.OBJECTFILE-01}), store the third variant with an
      artificially earlier \texttt{created} timestamp than its primary's
      own, and the fourth with one at or after it; construct the backend;
      confirm every variant is admitted, with the same size a directly
      stored object would receive, and only the damaged primary is
      reclaimed.
    \item With a test hook pausing the scan between classifying a damaged
      record and unlinking it, publish a fresh record at that path through
      a live store; resume the scan. Then, with a metadata update paused
      between its two writes (\reqref{FDS-STORAGE.METAUPDATE-01}), run the
      scan over that record and let the update finish.
    \item Read the reclaim counters once the backend is built; then read
      them and the scan's log summary after the scan.
    \item Repeat steps 1 and 3 in a build without the backend's
      \texttt{metrics} feature; read the log and the registry.
    \item Make the backend read-only, run the scan, and read its log summary.
    \item Make one shard directory unlistable and one record unreadable, run
      the scan, and read what it returns and logs.
  \end{enumerate}

  \textbf{Expected result:} A temporary is reclaimed exactly when its writer
  is not running, and a running writer's is never touched; a stopped
  writer's lease is removed and a running writer's kept; damaged, foreign,
  misplaced and withdrawn records are reclaimed; every variant is admitted, whatever
  its primary holds; no reclaimed entry, of any class, is ever admitted into the
  eviction manager; the fresh record published during the pause survives the
  resumed scan; the record read in the middle of a metadata update is
  admitted once the update finishes, not reclaimed; once the backend is
  built the counters read zero for every class, and after the scan they and
  the log summary report each class's reclaimed files and bytes; without
  the \texttt{metrics} feature the log summary is still written while no
  reclaim counter is registered; a scan
  of a read-only backend says in its summary that it stopped; the
  unlistable directory and the unreadable record are each skipped, named in
  a warning, and counted in what the scan returns and in its summary, and
  every other record is admitted;
  the scan requires no second pass over the tree beyond
  \reqref{TC-EVICTION.RECONSTRUCT-01}'s own.

  \reqdiaglink{overview}{\textit{View the domain overview.}}

  \reqcover{satisfy}{IMP}

  \verify{FDS-STORAGE.RECLAIM-01}
\end{requirement}

% ---------------------------------------------------------------------
% Directory sharding: shard-path derivation.
% ---------------------------------------------------------------------
\begin{requirement}{TC-LAYOUT.SHARDING-01}[draft][verify]{Shard layout and directory bound}

  \localreqdiagram*{tc22}{\reqid}{\reqdesc}

  \textbf{Objective:} The sharded path derived from an object's full
  \texttt{CacheKey} (used on lookup and store) is identical to the path
  derived from its \texttt{CompactCacheKey} (used on purge), for keys with and
  without a variance; the two-level split always falls on hex characters
  one-two and three-four of the object's name; the layout keeps every
  directory within the bound \reqref{REQ-EFFICIENCY.DIRSCALE-01} sets; and
  keys that differ only in their user tag are separate objects.

  \textbf{Setup:} The backend driven through the \texttt{Storage} interface,
  with a test hook that lets a test choose the name a key resolves to (a real
  hash cannot be steered to a chosen value).

  \textbf{Steps:}
  \begin{enumerate}
    \item Store an object under its \texttt{CacheKey}; record the shard
      path the store resolves to.
    \item Purge the same object using only its \texttt{CompactCacheKey};
      confirm the purge resolves to the identical shard path recorded in
      step~1.
    \item Repeat steps 1--2 for a variant: a key carrying a variance, purged
      through its own varianced \texttt{CompactCacheKey}.
    \item Through the hash hook, repeat with combined hashes that probe the
      split boundary (e.g.\ a hash whose first four characters repeat, and one
      where characters two and three differ sharply); confirm the first-level
      directory is always exactly characters one-two and the second-level
      directory always exactly characters three-four, for both key forms.
    \item Analysis: from the layout's fan-out, compute the expected and the
      high-percentile maximum entries per leaf directory for $10^8$ uniformly
      hashed objects, adding the temporaries a leaf can hold at peak
      concurrency.
    \item Store a large sample of objects with real hashes (e.g.\ $10^6$);
      count the entries in every first- and second-level directory and compare
      the distribution per leaf with the analysis scaled to the sample.
    \item Store two objects under keys for the same URL that differ only in
      their user tag; look both up; run the restart-time scan over the cache
      directory; then purge one and look both up again.
  \end{enumerate}

  \textbf{Expected result:} The path derived from the \texttt{CacheKey} and
  from the \texttt{CompactCacheKey} never diverge for the same object, with or
  without a variance; the shard boundary is exactly hex characters one-two
  (first level) and three-four (second level); the analysis shows fewer than
  2,000 entries in every directory at $10^8$ objects; the sample shows at most
  256 entries in every first-level directory and a per-leaf distribution
  consistent with the analysis; in step 7 each tagged key has a file of its
  own, each lookup returns its own body, the scan admits one eviction entry
  for each key, and after the purge only the purged key is a miss.

  \reqdiaglink{overview}{\textit{View the domain overview.}}

  \reqcover{satisfy}{IMP}

  \verify{SDS-LAYOUT.SHARDING-01}
  \verify{REQ-EFFICIENCY.DIRSCALE-01}
\end{requirement}

% ---------------------------------------------------------------------
% Shard-directory cache: startup rebuild and on-demand growth.
% ---------------------------------------------------------------------
\begin{requirement}{TC-LAYOUT.DIRCACHE-01}[draft][verify]{Shard cache matches disk at start}

  \localreqdiagram*{tc23}{\reqid}{\reqdesc}

  \textbf{Objective:} After construction, all 256 first-level directories
  exist and the in-memory cache exactly matches the second-level
  directories actually present on disk; concurrent first-time stores into
  the same new shard all succeed and leave it cached exactly once.

  \textbf{Setup:} The backend driven directly through the \texttt{Storage}
  interface by a test harness; a ``restart'' constructs a new backend instance
  over the same cache directory.

  \textbf{Steps:}
  \begin{enumerate}
    \item Pre-populate a subset of second-level directories directly on
      disk (bypassing the backend), then construct the backend; confirm
      all 256 first-level directories exist and the cache contains exactly
      the pre-populated second-level directories --- no more, no fewer.
    \item From multiple threads concurrently, store objects that resolve
      to the same, previously-unseen shard; confirm every store succeeds,
      the directory is created exactly once, and the cache ends up with
      exactly one entry for that shard.
    \item Restart the backend (a fresh construction) with objects now
      present under previously-created shards; confirm the newly built
      cache reflects them and that no second-level directory creation is
      attempted for them (first-level directories are always ensured at
      construction, per \reqref{FDS-LAYOUT.STARTUP-01}).
  \end{enumerate}

  \textbf{Expected result:} The in-memory cache never diverges from what
  is actually on disk; concurrent first-time stores into the same shard
  never fail or double-create; a restart always rebuilds an accurate cache
  from scratch.

  \reqdiaglink{overview}{\textit{View the domain overview.}}

  \reqcover{satisfy}{IMP}

  \verify{FDS-LAYOUT.STARTUP-01}
\end{requirement}

% ---------------------------------------------------------------------
% Object-file record layout: damage and truncation detection.
% ---------------------------------------------------------------------
\begin{requirement}{TC-LAYOUT.OBJECTFILE-01}[draft][verify]{Damaged records treated as a miss}

  \localreqdiagram*{tc24}{\reqid}{\reqdesc}

  \textbf{Objective:} A reader accepts a published record only when its
  magic, version, header and metadata CRCs, and stored key all validate, its
  header's lengths fit the record and the file, and its length is exactly
  the computed body offset plus \texttt{body\_len}; any single violation is a miss that leaves the file in
  place; a corrupted body is read to a normal end, and logged when the
  read-time body check is on; and a metadata read racing an in-place update
  sees one complete version.

  \textbf{Setup:} The backend driven directly through the \texttt{Storage}
  interface by a test harness; a ``restart'' constructs a new backend instance
  over the same cache directory. The steps run in the default build and again
  in a build with the \texttt{body-crc-verify} feature.

  \textbf{Steps:}
  \begin{enumerate}
    \item Store an object normally, then flip a byte in its on-disk
      \texttt{magic} or \texttt{version} field; look it up.
    \item Store an object, then flip a byte within the header or metadata
      region without updating the corresponding CRC; look it up.
    \item Store an object, then overwrite its \texttt{version} with another
      value and a matching \texttt{header\_crc}, as a record written by
      another format revision would carry; look it up.
    \item Store an object, then copy its file over another key's canonical
      path, so the stored key differs from the looked-up one; look up that
      other key.
    \item Store an object, then truncate its file short of, and separately
      extend it past, the computed body offset plus \texttt{body\_len}; look
      it up.
    \item Store an object whose body spans several read chunks, then flip a
      byte within its body without changing its length; look it up and read
      the whole body. In the \texttt{body-crc-verify} build, repeat with the
      read-time body check switched on, and with it on also read an intact
      body whole and the corrupted body from a sought position.
    \item Repeatedly apply in-place metadata updates to one object while
      concurrent lookups read it.
    \item Store an object, then rewrite one of its header's lengths with a
      matching \texttt{header\_crc}: \texttt{meta\_capacity} so small that
      the metadata segments no longer fit, and, separately,
      \texttt{key\_len} so large that the key and the body would lie past
      the end of the file; look it up after each.
  \end{enumerate}

  \textbf{Expected result:} Each lookup in steps 1--5 and 8 reports a miss, and the
  damaged file is still present afterwards; every body read in step 6 ends
  normally, and the backend logs an error naming the file for the whole read
  of the corrupted body with the check on, and for no other read; every
  lookup in step 7 returns
  either the metadata before an update or after it, never a mixture and never
  a miss while no update is running. A store that never
  reaches completion is covered separately
  (\reqref{TC-COMPLIANCE.INCOMPLETE-01}): it is never visible under the
  canonical name at all, so this test case's scope is damage to an
  already-published record, not incompleteness.

  \reqdiaglink{overview}{\textit{View the domain overview.}}

  \reqcover{satisfy}{IMP}

  \verify{SDS-LAYOUT.OBJECTFILE-01}
\end{requirement}

\begin{requirement}{TC-LAYOUT.KEYSEGMENT-01}[draft][verify]{Key segment round trip}

  \localreqdiagram*{tc36}{\reqid}{\reqdesc}

  \textbf{Objective:} A key is stored in exactly the encoding
  \reqref{SDS-LAYOUT.KEYSEGMENT-01} specifies and decodes back to the same
  key; a key segment that does not follow the encoding is rejected, and the
  restart-time scan reclaims its record.

  \textbf{Setup:} For steps 1 and 2, the key segment's encoder and decoder
  driven directly by unit tests. For step 3, the backend driven through the
  \texttt{Storage} interface, and its restart-time scan run over the same
  cache directory.

  \textbf{Steps:}
  \begin{enumerate}
    \item Encode keys with and without a variance hash, each with an empty
      and a non-empty tag; inspect each encoding, then decode it.
    \item Decode a segment shorter than the primary hash and the variance
      flag; a segment whose flag is 1 but which ends before the variance hash
      does; a segment whose flag is 2; and a segment whose tag is not valid
      UTF-8.
    \item Store an object, then set its stored key's variance flag to 2 and
      reseal the record's metadata and header checksums over the change; run
      the restart-time scan.
  \end{enumerate}

  \textbf{Expected result:} In step 1 every field lies at the offset the
  encoding specifies, the tag runs to the end of the segment, and every key
  decodes to itself; every decode in step 2 fails; in step 3 the scan
  reclaims the record as damaged and does not admit it.

  \reqdiaglink{overview}{\textit{View the domain overview.}}

  \reqcover{satisfy}{IMP}

  \verify{SDS-LAYOUT.KEYSEGMENT-01}
\end{requirement}
